\subsection{复向量空间上的弱拓扑}

II, § 6 No. 1 与 No. 2 的定义和结果可不加改变地适用于复向量空间。若 F 与 G 是由双线性型 B 处于对偶的两个复向量空间，则底空间 \(F_{0}\) 与 \(G_{0}\) 由 RB 处于对偶，且由 II, p.~61 公式 (1) 推出，弱拓扑 \(\sigma(\mathrm{F},\mathrm{G})\) 与 \(\sigma(\mathrm{F}_{0},\mathrm{G}_{0})\) 是相同的。

定义 1. --- 设 F 与 G 是处于对偶的两个复向量空间。对 F 的任意子集 M，M 在 G 中的极集，记作 \(M^{\circ}\)，是 \(y \in G\) 中使 \(\mathcal{R}(\langle x, y \rangle) \geqslant -1\)（对所有 \(x \in M\)）者所成的集。

若 \(\mathbf{M}^{\circ}\) 是 \(\mathbf{M} \subset \mathbf{F}\) 在 \(\mathbf{G}\) 中的极集，则 \((\lambda \mathbf{M})^{\circ} = \lambda^{-1}\mathbf{M}^{\circ}\) 对所有 \(\lambda \in \mathbf{C}^{*}\) 成立。

若 M 是 E 的一个（复）向量子空间，则 \(M^{\circ}\) 是一个闭向量子空间（对 \(\sigma(\mathbf{G},\mathbf{F})\) 而言），因为关系 \(\mathcal{R}(\lambda\langle x,y\rangle)\geqslant-1\)（对每个标量 \(\lambda\in C\)）蕴含 \(\langle x,y\rangle=0\)；我们仍说 \(M^{\circ}\) 是 G 中与 M 正交的子空间。

若 M 是 F 中的一个均衡集，则 \(M^{\circ}\) 是 G 中的一个均衡集；这时 \(M^{\circ}\) 是 \(y \in G\) 中使 \(\left|\langle x, y \rangle\right| \leqslant 1\)（对所有 \(x \in M\)）者所成的集；因为这个关系等价于 \(\mathcal{R}(\langle \zeta x, y \rangle) \leqslant 1\) 对所有 \(x \in M\) 与所有 \(\zeta \in C\)（满足 \(|\zeta| = 1\)）成立。

II, p.~41 至 51 的结果对复向量空间也无限制地成立。

